\section{Tau-bench Coding Walkthrough：从论文接口到 AgentBeats}

课程最后用 Tau-bench 展示如何把论文 benchmark 转为多 Agent 服务。实现分为 interface analysis、workflow design、kickoff script、Green Agent、White Agent 与平台 integration。重点不是复制代码截图，而是理解每个进程拥有的状态和协议。

\subsection{先梳理 interface：谁知道答案，谁可以调用工具}

slide 91 给出原则：被测 Agent 应能在给定 task 上求解；Green Agent 提供允许的工具并像真实用户一样交互。slide 92 要求先读论文，再读 codebase，写出 task formulation、user simulator、tools、database 和 verifier 的映射。slide 94 的最终 build 指出两个挑战：cross-agent tool use 与迁移原 evaluation。

\lecturefigure{slides-images/slide-091.png}{官方 slide 91：Interface 原则与 Green/White 分工}
\lecturefigure{slides-images/slide-092.png}{官方 slide 92：从论文 task formulation 到 codebase}
\lecturefigure{slides-images/slide-094.png}{官方 slide 94：跨 Agent tool use 与 evaluation migration 两个挑战}

\begin{importantbox}{Interface inventory}
列出每个角色的输入、输出、可调用工具、持有状态和秘密：White Agent 持对话与公开工具；user simulator 持隐藏用户目标；environment 持数据库；Green Agent 持 task/metric 与 orchestration；平台持 run metadata。任何 secret 同时出现在 participant prompt 都会造成泄漏。
\end{importantbox}

\subsection{Workflow：Kickoff、对话、终止与评分}

slide 95 给出 workflow：kickoff script 连接平台与 agents，发送启动消息；Green Agent 开始测试并路由交互；结束后调用 metric，提交 report。需要定义谁先发言、turn limit、tool routing、用户模拟结束条件、White Agent 提交 final 的格式，以及异常如何编码。

\lecturefigure{slides-images/slide-095.png}{官方 slide 95：Kickoff script、Green Agent 与 White Agent workflow}

\begin{knowledgebox}{消息 envelope}
建议每条消息含 run\_id、sender、recipient、type、payload、timestamp、sequence 与 trace\_id。Tool result 使用结构化 schema，不把异常堆栈直接拼入自然语言。sequence 防止乱序，run\_id 防止并发串线，trace\_id 连接平台日志与环境动作。
\end{knowledgebox}

\subsection{Kickoff script：最小控制平面}

slide 96 的代码展示启动入口。Kickoff 解析 agent card/endpoint，创建或连接 run，向 Green Agent 发送测试请求并等待结果。它不应包含 benchmark 评分逻辑，否则平台调用和 benchmark 语义耦合。重试必须幂等，避免网络超时后重复创建环境。

\lecturefigure{slides-images/slide-096.png}{官方 slide 96：Kickoff script 代码结构}

\begin{warningbox}{启动脚本常见故障}
硬编码 endpoint、没有 timeout、把 credential 写日志、重试重复计分、进程退出后环境未销毁。应使用配置/secret manager、明确 deadline、idempotency key 和 finally cleanup。脚本越小，越容易在不同平台重用。
\end{warningbox}

\subsection{Green Agent：Orchestrator、User Simulator 与 Evaluator}

slide 97 展示 Green Agent 实现。它加载 task，初始化环境，与 White Agent 多轮交互，必要时扮演用户或调用 user simulator，最后读取 final state 并评分。自然语言交互与确定性 metric 应分层：LLM 可生成用户行为，代码负责权限、状态和硬指标。

\lecturefigure{slides-images/slide-097.png}{官方 slide 97：Green Agent 的 orchestrator 与 evaluator 实现}

\begin{knowledgebox}{Green Agent 内部状态}
至少保存 task spec、environment handle、turn count、messages、tool trace、budget、termination reason、raw metric evidence 与 final report。不要只保存聊天文本；数据库 diff、API response 和 verifier output 是 outcome validity 的关键证据。
\end{knowledgebox}

\subsection{White Agent：被测策略必须可替换}

slide 98 展示 White Agent 作为 participant。它接收用户/Green 消息，调用公开工具，生成回复。benchmark 不应依赖某个特定模型 SDK；用 Agent Card/A2A 或稳定接口描述 capability，使不同实现都能被同一 Green Agent 测试。White Agent 的本地日志可用于调试，但正式评分只依赖 Green/环境侧证据。

\lecturefigure{slides-images/slide-098.png}{官方 slide 98：White Agent participant 实现}

\begin{warningbox}{不要让 White Agent 自报成功}
participant 的 ``done'' 只表示请求结束，不是任务成功。Green Agent 必须独立读取环境状态并运行 verifier。若让 White Agent返回 score 或 ground truth comparison，模型可以直接优化自我声明，评估失去可信性。
\end{warningbox}

\subsection{AgentBeats integration：可发现、可复现、可观测}

slide 99 列出实现完成后的平台工作：hosting、access、reproducible/open、report result/trace、package protocol。slide 102 展示 Google ADK 等开发/日志工具。集成的目标是让平台发现 agent、启动相同版本、查看 trace 并获得结构化 metric，而不是要求所有项目使用同一内部框架。

\lecturefigure{slides-images/slide-099.png}{官方 slide 99：AgentBeats hosting、访问、复现、trace 与协议集成}
\lecturefigure{slides-images/slide-102.png}{官方 slide 102：ADK 与在线 logging 工具}

\begin{importantbox}{可复现发布包}
包含源码 commit、容器镜像 digest、依赖锁、agent card、环境变量 schema、任务数据版本、运行命令、sample trace、baseline result 和许可证。秘密与私有测试通过受控存储注入，不进入镜像。平台结果应引用所有 digest，保证半年后能重跑。
\end{importantbox}

\begin{knowledgebox}{可观测性最小集合}
run timeline、消息 trace、tool latency/error、token/compute cost、environment diff、verifier evidence、termination reason。公开 leaderboard 可只展示汇总，但项目调试和审计必须保留细粒度证据，并对敏感数据做脱敏和 retention policy。
\end{knowledgebox}

\subsection{本章小结}

Tau-bench walkthrough 展示了从论文到评估服务的完整迁移：先冻结角色与 secret，设计协议和状态机，再分别实现 kickoff、Green、White 与环境 verifier，最后用可复现包接入平台。接口正确性和 evidence chain 比“代码看起来 Agentic”更重要。

